\part{The pipeline blocks}

% =====================================================================
\chapter{What these four blocks replace}
\label{ch:pipeline}

Everything up to here has assumed that the files a run needs already exist: a
\file{geometry.dat} somebody made, a metabolic model somebody converted, a
surrogate somebody trained, and a habit of loading the output into Python
afterwards to find out whether anything was conserved.

Those four assumptions were four scripts, run by hand, in the right order,
before and after every simulation. This part is about the four XML blocks that
do them instead.

\begin{center}
\begin{tikzpicture}[font=\footnotesize,
  b/.style={draw=cbRule,fill=cbSand,rounded corners=2pt,minimum width=4.3cm,
            minimum height=0.85cm,align=center,line width=0.7pt},
  o/.style={draw=cbRule,fill=white,rounded corners=2pt,minimum width=4.3cm,
            minimum height=0.85cm,align=center,line width=0.7pt},
  a/.style={-{Latex[length=2mm]},cbGrey,line width=0.7pt}]

\node[o] (g1) at (-3.4, 1.4) {\texttt{geometry.py}\\ by hand, before};
\node[o] (g2) at (-3.4, 0.0) {\texttt{extractMM.py}\\ by hand, before};
\node[o] (g3) at (-3.4,-1.4) {four training scripts\\ by hand, before};
\node[o] (g4) at (-3.4,-2.8) {\texttt{postprocess.py}\\ by hand, after};

\node[b] (n1) at ( 3.4, 1.4) {\tg{domain}\tg{generate}};
\node[b] (n2) at ( 3.4, 0.0) {\tg{model\_source}};
\node[b] (n3) at ( 3.4,-1.4) {\tg{surrogate}};
\node[b] (n4) at ( 3.4,-2.8) {\tg{diagnostics}};

\draw[a] (g1) -- (n1);
\draw[a] (g2) -- (n2);
\draw[a] (g3) -- (n3);
\draw[a] (g4) -- (n4);
\end{tikzpicture}
\end{center}

\note{The old way still works.}{None of these blocks is required, and every
one of them is absent by default. An input file written before they existed
takes exactly the path it always took --- the code checks for the block, does
not find it, and carries on. The Python tools in \file{tools/} are still
shipped and still supported; they do things C++ cannot, such as reading a TIFF
stack or drawing a plot.}

\warn{One thing genuinely changed.}{\tg{model\_filename} may now point at an
SBML file directly, and \file{extractMM.py} is no longer needed for the GLPK
path. The old matrix files still load --- the format is decided by looking at
the file's root element, not its name --- so nothing you already have stops
working.}

% =====================================================================
\chapter{Where your metabolic model comes from}
\label{ch:modelsource}

\section{The point}

A genome-scale model is a published artefact with a revision. Chapter
\ref{ch:reference} treats \tg{model\_filename} as just a path, but in practice
getting that path right was three manual steps: find the model on BiGG,
download it, run \file{extractMM.py} to flatten it into the matrix format
CompLaB used to read. Each step is somewhere the wrong file can be substituted
without anything complaining.

\section{Reading SBML directly}

The simplest change is that \tg{model\_filename} can now name the SBML file
itself:

\begin{lstlisting}[style=xml]
<model_filename>iJO1366.xml</model_filename>
\end{lstlisting}

The reader is \file{complab3d\_sbml.hh}. It produces the same
$S$, $b$, $c$, $lb$, $ub$ and objective column that \file{extractMM.py} did, in
the same order, and was checked against \texttt{cobrapy} on \emph{iJO1366}:
identical reaction order, identical $4.7$ million-entry stoichiometry matrix,
identical bounds, identical objective. It takes about $0.35$\,s where
\texttt{cobrapy} takes $5$\,s, because it does one pass of tinyxml and no
object graph.

\note{How the format is decided.}{By the root element. \code{<sbml>} means
SBML; \code{<Metabolic\_Model>} means the \file{extractMM.py} matrix format.
The file extension is never consulted, so renaming a file cannot change how it
is read.}

\section{Naming exchange reactions instead of numbering them}

This is the part worth changing in an existing input file even if you change
nothing else.

\begin{lstlisting}[style=xml]
<!-- before: column numbers, correct only for one revision of one model -->
<exchange_reaction_indices>27 35 -1</exchange_reaction_indices>

<!-- after: names, checked against the model at start-up -->
<exchange_reaction_names>EX_glc__D_e EX_o2_e none</exchange_reaction_names>
\end{lstlisting}

Both mean the same thing --- which reaction of the metabolic model exchanges
each substrate, in \tg{name\_of\_substrates} order. The difference is what
happens when they are wrong.

\begin{itemize}[leftmargin=1.4em]
\item An index that points at the wrong reaction is not detectable. The run
      proceeds, the linear program is well posed, and the answer is wrong
      everywhere.
\item A name that does not exist stops the run at start-up, names the
      substrate, and lists the boundary reactions whose names are close.
\end{itemize}

Use \code{none} (or \code{-1}) for a substrate this organism does not exchange.
If both tags are given, the names win and a note says so.

\warn{Names need SBML.}{The \file{extractMM.py} matrix format does not carry
reaction names, so \tg{exchange\_reaction\_names} with a matrix file is a hard
error rather than a silent fallback. Point \tg{model\_filename} at the SBML.}

\section{Letting the run find the model}

\begin{lstlisting}[style=xml]
<model_source>bigg:e_coli_core</model_source>
<model_bundle>models</model_bundle>      <!-- shipped with the code -->
<model_cache>input</model_cache>         <!-- where the unpacked file goes -->
<allow_download>false</allow_download>   <!-- default: never touch the network -->
\end{lstlisting}

The search order is fixed and is reported in the log:

\begin{enumerate}[leftmargin=1.6em]
\item \textbf{already unpacked} in \tg{model\_cache}. Used as is.
\item \textbf{bundled} with the code as \file{models/<id>.xml.gz}. Unpacked
      into the cache. Three models ship this way: \emph{e\_coli\_core},
      \emph{iJO1366} and \emph{STM\_v1\_0}.
\item \textbf{downloaded}, and only if \tg{allow\_download} is \code{true}.
      A cluster compute node usually has no route to the internet, which is
      exactly why the default is \code{false} and why the models are bundled.
\end{enumerate}

Whichever branch runs, the file is checked against \file{models/manifest.txt},
which records for each model its species count, reaction count, parsed
metabolite and reaction counts, objective reaction, SHA-style hash, byte size
and source URL. A mismatch does not stop the run; it prints what it expected
and what it found, so a model that has been revised upstream is visible in the
log rather than buried in the results.

\note{One microbe, or all of them.}{The resolved file is given to every
metabolic microbe that did not name its own \tg{model\_filename}. A microbe
with an explicit one keeps it, so a two-organism run can mix a bundled model
with a local one.}

\note{Parallel runs.}{Only rank 0 unpacks or downloads. The others wait at a
barrier and then find the file already there. Nothing is written twice and no
two ranks race for the same path.}

% =====================================================================
\chapter{Making a pore space}
\label{ch:geomgen}

\section{The point}

\file{tools/geometry.py} can build a packing or threshold an image stack, and
it is still the right tool when the source is a TIFF stack or a PNG series.
But for the common cases --- a sphere pack at a given porosity, a fracture, a
layered medium --- there is no reason to leave the input file.

\section{Generating}

Inside \tg{domain}, alongside \tg{nx} and friends:

\begin{lstlisting}[style=xml]
<generate>spheres</generate>      <!-- channel | spheres | cylinders |
                                       random | fracture | layered -->
<porosity>0.55</porosity>
<grain_radius>3.0</grain_radius>
<seed>12345</seed>
<walls>y</walls>                  <!-- solid walls on the named axes -->
<write_geometry>generated.dat</write_geometry>
\end{lstlisting}

The generated volume is written to \tg{write\_geometry} inside the input
directory before the run starts. That file is the record: the run is
reproducible from its own output even if the seed, the code or the XML later
change, and the generated geometry can be inspected with exactly the same
tools as one that came from imaging.

\section{Importing a binary volume}

\begin{lstlisting}[style=xml]
<import_raw>rock.raw</import_raw>
<raw_dtype>uint8</raw_dtype>
<threshold>128</threshold>
<invert>false</invert>
\end{lstlisting}

A flat binary volume, thresholded into pore and solid. \tg{generate} and
\tg{import\_raw} are mutually exclusive and setting both stops the run.

\section{What the inspection tells you}

Either way, the pore space is inspected before the flow solver is built, and
the report is printed:

\begin{lstlisting}[style=sh]
[GEOM] generated a 'spheres' pore space
[GEOM] 24 x 24 x 6 = 3456 voxels
[GEOM]   material 1   bounce_back    288    8.33%
[GEOM]   material 0   solid         1170   33.85%
[GEOM]   material 2   pore          1998   57.81%
[GEOM] porosity 0.5781
[GEOM] percolates along x: yes, 1 spanning cluster(s)
[GEOM] isolated pore space: 1 voxel(s), 0.1% of the open space
[GEOM] written to input/generated.dat so this run can be reproduced exactly
\end{lstlisting}

\warn{A sealed domain stops the run.}{If the pore space does not percolate
along $x$, the run stops before the flow solver starts. A pressure drop across
a sealed medium has no solution, and the flow solver would spend its entire
iteration budget failing to converge to one. Lower the porosity target, or
shrink the grains.}

\note{Isolated pockets are reported, not removed.}{Pore voxels that no
percolating path reaches still hold substrate and still host biology, and
nothing about them is wrong --- but they never exchange with the through-flow,
so a mass balance that includes them behaves differently from one that does
not. The count is printed so the number is not a surprise.}

\note{The file is \code{nx} slices wide, where \code{nx} is what you wrote in
the XML.}{Internally the code adds two ghost columns, and \code{readGeometry}
duplicates the first and last slice into them. The generator is given the
width you asked for, so this is handled --- but it is why a geometry file made
by hand must have exactly $\code{nx}\times\code{ny}\times\code{nz}$ entries and
not two slices more.}

% =====================================================================
\chapter{A surrogate without leaving the XML}
\label{ch:surrogateblock}

\section{The point}

Chapter \ref{ch:trainsurrogate} describes the offline route: sweep the FBA into
a CSV, fit a network, export it as a header, paste the header into
\file{surrogateModel.hh}, rebuild. That route still exists and is still the
right one when you want to inspect the fit, try several architectures, or
train on a machine that is not the one you will run on.

The change is that the weights now live in a \textbf{file the solver reads at
run time} rather than in a header it was compiled with. Once that is true, the
run can also fit the network itself.

\section{The block}

\begin{lstlisting}[style=xml]
<surrogate>
    <enabled>true</enabled>
    <weights_file>output/ecoli.srg</weights_file>
    <train_if_missing>true</train_if_missing>
    <train>
        <samples>2000</samples>
        <layers>10 10</layers>
        <restarts>5</restarts>
        <epochs>4000</epochs>
        <verify_points>512</verify_points>
        <seed>20260814</seed>
        <input0>S 0.1 10.0</input0>
        <input1>O 0.1 20.0 log</input1>
    </train>
</surrogate>
\end{lstlisting}

Each \tg{inputN} line reads: \emph{substrate name}, \emph{low}, \emph{high},
and optionally \code{log} for a logarithmic sweep. The substrate name must be
one of \tg{name\_of\_substrates}.

\section{What happens at start-up}

\begin{center}
\begin{tikzpicture}[font=\footnotesize,
  s/.style={draw=cbRule,fill=cbSand,rounded corners=2pt,minimum width=6.4cm,
            minimum height=0.85cm,align=center,line width=0.7pt},
  a/.style={-{Latex[length=2mm]},cbGrey,line width=0.7pt}]
\node[s] (a1) at (0, 0)    {\tg{weights\_file} exists?};
\node[s] (a2) at (0,-1.5)  {load it, print its valid range};
\node[s] (b2) at (0,-3.2)  {build a linear program from the model};
\node[s] (b3) at (0,-4.7)  {sweep it, fit, verify on fresh points};
\node[s] (b4) at (0,-6.2)  {write the \texttt{.srg}; later runs load it};
\draw[a] (a1) -- node[right,font=\tiny,xshift=2pt]{yes} (a2);
\draw[a] (a1.west) -- ++(-2.6,0) |- node[left,font=\tiny,pos=0.25]{no, and \tg{train\_if\_missing}} (b2.west);
\draw[a] (b2) -- (b3);
\draw[a] (b3) -- (b4);
\end{tikzpicture}
\end{center}

The linear program the sweep solves is not a copy. If the run already has a
GLPK microbe, its own persistent \code{glp\_prob} is used --- the same object,
with the same \tg{constraint\_indices} overrides, that the simulation would
have solved. If it does not (the usual case: one organism, reaction type
\code{surrogate}), a temporary problem is built from that microbe's own model
and released as soon as training finishes.

\note{Positive uptake in, negative bound out.}{Exchange reactions are written
\code{metabolite ->}, so consuming is a negative flux. The sweep passes
positive uptake rates in \code{mmol/gDW/h}, the same convention as
\code{Fin} everywhere else, and the sign is applied in exactly one place. An
infeasible solve is recorded as zero growth rather than an error, so the
network learns where the feasible region ends instead of having a hole in its
domain.}

\section{Training is on rank 0 only}

Every rank fitting its own network from its own random start would give a
different growth rate in each subdomain. Rank 0 trains, writes the file,
everyone barriers, the others read it.

\section{Which substrate is which}

A \file{.srg} file records the substrate name of every input:

\begin{lstlisting}[style=sh]
inputs 2
inputnames S O
trainmin 0.127038797 0.118747061
trainmax 9.983995766 19.978325403
\end{lstlisting}

At start-up those names are matched against \tg{name\_of\_substrates}, and a
name that is not there stops the run. Without this a network trained on
(acetate, Fe\textsuperscript{3+}) could be fed them the other way round: every
number plausible, every growth rate wrong, nothing to see in the log.

\warn{A file with no \code{inputnames} line.}{Networks exported by the older
MATLAB route have no such line. They still load, and their inputs are assumed
to be the first substrates in order --- with a warning saying so. Check it.}

\section{Outside the training box}

A neural network does not fail outside the range it was fitted on. It returns
a confident number that means nothing. So the range travels with the weights,
and at run time every evaluation is \textbf{clamped} to that box and
\textbf{counted}:

\begin{lstlisting}[style=sh]
[SRG] runtime network: 1076692 evaluation(s), 0 clamped (0.00%)
\end{lstlisting}

Clamping rather than extrapolating is the conservative choice: at the edge the
network returns the growth rate at the nearest uptake that was actually solved
for, which is a real answer for a nearby state. It is still an approximation,
which is why the count is reported. Above 5\% the report says, in as many
words, that the results should not be used until the ranges are widened and
the network refitted.

\section{Does it agree with the FBA?}

On \emph{e\_coli\_core}, in a run otherwise identical, glucose and oxygen as
the two inputs:

\begin{center}
\begin{tabular}{@{}lrr@{}}
\toprule
& \textbf{GLPK, every voxel} & \textbf{surrogate} \\
\midrule
biomass change over the run & $-4.3090\%$ & $-4.3107\%$ \\
final product concentration & $2.0871\times10^{-2}$ & $2.0871\times10^{-2}$ \\
wall clock                  & $492$\,s & $3.3$\,s \\
\bottomrule
\end{tabular}
\end{center}

A factor of about $150$ in wall clock, for a difference in the answer of four
parts in ten thousand. That ratio is the whole argument for surrogates, and it
is why the training range matters so much: the speed is only worth having if
the simulation stays inside the box.

% =====================================================================
\chapter{The run's own record}
\label{ch:diagnostics}

\section{The point}

Until this block existed, a CompLB3D run wrote VTI volumes and nothing else.
Every scalar question --- did mass balance, how did porosity change, when did a
concentration first go negative --- meant loading the volumes into Python
afterwards and hoping the answer was still in them.

\section{The block}

\begin{lstlisting}[style=xml]
<diagnostics>
    <enabled>true</enabled>
    <summary_csv>summary.csv</summary_csv>
    <interval>100</interval>            <!-- 0 = follow the VTI interval -->
    <tolerance>1e-6</tolerance>
    <conserve>Fe2 + Fe3</conserve>
    <conserve1>SO4 + HS</conserve1>
</diagnostics>
\end{lstlisting}

\section{summary.csv}

One row per interval, written next to the VTI files:

\begin{lstlisting}[style=sh]
iteration,porosity,open_voxels,S_total,S_mean,S_min,S_max,O_total,...
0,0.578125,1998,1998,1,0,1
100,0.578125,1998,1832.4,0.917,0,1
\end{lstlisting}

Totals and means are over \textbf{open voxels only} --- pore plus biofilm --- and
over the physical domain $x = 1 \ldots n_x-2$, excluding the two ghost columns.
Both choices matter: including the solid would make the mean fall as biofilm
grew for no chemical reason, and including the ghost columns would count the
inlet and outlet slices twice, so a conserved total would jump whenever the
inlet concentration changed.

\section{The conservation checks}

\tg{conserve} names a sum of substrates that ought to be constant. The first
recorded value is the baseline; after that the drift is
\[
  \text{drift} \;=\; \frac{|\Sigma_{\text{now}} - \Sigma_{\text{first}}|}
                          {\max\!\left(|\Sigma_{\text{first}}|,\,
                                        |\Sigma_{\text{now}}|\right)} .
\]
Normalising by the larger of the two, rather than by the baseline, is what
keeps the number meaningful when a run starts from an empty domain --- there
the baseline is zero and dividing by it reports $10^{302}$, which tells nobody
anything.

\warn{What to conserve, and what not to.}{A substrate held at a fixed
concentration on an inlet is being supplied from outside. Its total is not
conserved and never will be, and asking for it to be produces a failure that
is entirely your own. Conserve a \emph{closed} sum instead: every species
sharing one element, which is a conserved moiety of the reaction network. The
failure message says this, in those words, because it is by far the most
common cause.}

A negative concentration is reported once, loudly, the first time it appears,
with the substrate and the step. At the end of the run each check is reported
as PASS, FAIL, or SKIPPED:

\begin{lstlisting}[style=sh]
[DIAG] summary written to output/summary.csv
[DIAG] Fe2 + Fe3            worst relative drift 3.104e-09  PASS
[DIAG] SO4 + HS             worst relative drift 1.882e-02  FAIL
[DIAG] verdict: FAIL
\end{lstlisting}

\note{SKIPPED is not PASS.}{A \tg{conserve} naming a substrate that does not
exist is reported as SKIPPED, never as PASS. A skipped check that looks like a
pass is how a mass-balance guarantee quietly becomes worthless.}

\note{Something this found.}{Running example 02 with no-flux boundaries at both
ends --- a genuinely closed box --- the tracer total rises about
2.3\% over the first two hundred steps and is then flat to
$10^{-5}$ for the rest of the run. That is a start-up transient of the Neumann
boundary implementation, not an ongoing leak, and it is present with or without
any of the new code. It is recorded here because it is exactly the kind of
thing this block exists to surface, and because a tolerance below about
$3\times10^{-2}$ on a Neumann case will report it.}
